\section*{Capítulo \ref{ch:b2:boltzmann-factor} --- El factor de Boltzmann}

\begin{solution}{exo:b2:boltzmann-factor:1}
$H = k_BT/mg = \qty{8.4}{km}$; \qty{0.70}{}, \qty{0.35}{}, \qty{0.30}{bar} (los
\qty{0.26}{} reales: el aire se enfría con la altura); helio \qty{61}{km}; Marte
\qty{10.7}{km}.
\end{solution}

\begin{solution}{exo:b2:boltzmann-factor:2}
Hidrógeno: $4\eu^{-395} \approx 10^{-171}$; $4\eu^{-19.7} = 10^{-8}$; $4\eu^{-11.8} = 3 \times
10^{-5}$. Sodio: $\eu^{-9.8} = 6 \times 10^{-5}$. Rotación: $\eu^{-0.04} = 0.96$, casi
iguales. Espín nuclear: $1 - 4 \times 10^{-6}$ --- la polarización minúscula con la que
trabaja la resonancia magnética.
\end{solution}

\begin{solution}{exo:b2:boltzmann-factor:3}
N$_2$ a \qty{300}{K}: \qty{422}{}, \qty{476}{}, \qty{517}{m/s}; a \qty{77}{K}: \qty{214}{}, \qty{241}{},
\qty{262}{m/s}; H$_2$: \qty{1580}{}, \qty{1780}{}, \qty{1930}{m/s}; partícula de humo:
\qty{9}{cm/s} (su tiritona browniana). Por encima de $2v_{\text{p}}$, una molécula de
cada veinte; por encima de $3v_{\text{p}}$, una de cada $2500$ --- todavía $10^{22}$ por metro cúbico.
\end{solution}

\begin{solution}{exo:b2:boltzmann-factor:4}
$x = 2.2 \times 10^{-3}$ y $0.67$; $M/M_{\text{sat}} = 2.2 \times 10^{-3}$ y $0.59$;
$M_{\text{sat}} = n\mu = \qty{9.3e4}{A/m}$; $\chi = \mu_0n\mu^2/k_BT = 2.6 \times 10^{-4}$.
Los \qty{1.7e6}{A/m} del hierro son cooperativos, no un efecto de Boltzmann.
\end{solution}

\begin{solution}{exo:b2:boltzmann-factor:5}
(a) \cref{prop:b2:boltzmann-factor:twolevel}. (b) $\eu^{-x} < 1$ para $T > 0$;
$N_2 > N_1$ exigiría $T < 0$ --- una \omterm{prop:b2:laser:gain}{inversión de población}. (c) $0$ y
$Nk_B\ln 2$. (d) $x = 1.93$: $0.13$; a \qty{77}{K}, $x = 7.5$: $5 \times 10^{-4}$.
\end{solution}

\begin{solution}{exo:b2:boltzmann-factor:6}
(a) Derivada de $N\varepsilon/(\eu^x + 1)$. (b) $x = 2.40$, $C = 0.44\,Nk_B$.
(c) $\propto x^2\eu^{-x} \to 0$: congelado; $\propto 1/T^2 \to 0$: ambos niveles llenos,
sin nada que absorber. (d) $k_BT = 0.42\varepsilon$: \qty{0.5}{mK}; por debajo de un
milikelvin los espines nucleares acaparan casi toda la capacidad calorífica y
frenan cualquier enfriamiento.
\end{solution}

\begin{solution}{exo:b2:boltzmann-factor:7}
(a) \cref{thm:b2:boltzmann-factor:equipartition}. (b) $\tfrac32R$ y
$\tfrac52R$. (c) $C_{\text{vib}}/R = x^2\eu^x/(\eu^x - 1)^2$: $0.002$, $0.41$, $0.90$;
$C_V = 2.50R$, $2.91R$, $3.40R$. (d) $0.92$, $0.30$, $0.995$; diamante con $x = 4.3$:
$0.26 \times 3R = \qty{6.5}{J/mol/K}$.
\end{solution}

\begin{solution}{exo:b2:boltzmann-factor:8}
(a) y (b) \cref{thm:b2:boltzmann-factor:maxwell}. (c) $\tfrac32k_BT$. (d)
$\tfrac14n\langle v\rangle = \qty{2.9e27}{m^{-2}.s^{-1}}$: $3 \times 10^{23}$ por centímetro
cuadrado y segundo; por \qty{7.9e-13}{m^2}: $2.3 \times 10^{15}$
moléculas por segundo, una fuga de $\qty{1e-5}{Pa.m^3/s}$.
\end{solution}

\begin{solution}{exo:b2:boltzmann-factor:9}
(a) Intégrese $f$ por partes; el término dominante es $(2/\sqrt\pi)x\eu^{-x^2}$. (b)
Helio con $x = 5.5$: $5 \times 10^{-13}$; nitrógeno con $x = 14.5$: $\eu^{-211}$. (c) $x =
4.9$: $10^{-10}$ por tiempo de colisión --- desaparecido; $10^{-17}$ exige $x \approx 6.6$,
$T \lesssim \qty{230}{K}$. (d) $x = 11$: $\eu^{-121}$; bastante frío.
\end{solution}

\begin{solution}{exo:b2:boltzmann-factor:10}
(a) $n = n_0\eu^{m\omega^2r^2/2k_BT}$. (b) $\eu^{\Delta m\omega^2r^2/2k_BT} = \eu^{0.24} =
1.27$. (c) $H = k_BT/m'g = \qty{16}{\micro m}$; contar esferas a
varias alturas da $k_B$ y, con él, $N_A = R/k_B$. (d) Los iones siguen
$\eu^{\mp e\varphi/k_BT}$ alrededor de una carga; linealizado en la \omterm{prop:b2:maxwell-equations:potentials}{ecuación de Poisson}
esto da $\varphi'' = \varphi/\lambda_{\text{D}}^2$: el potencial queda apantallado
más allá de $\lambda_{\text{D}}$.
\end{solution}

\begin{solution}{exo:b2:boltzmann-factor:11}
(a) $Z = 1/(1 - \eu^{-x})$, $\langle E\rangle = h\nu/(\eu^x - 1)$. (b) $k_BT$:
equipartición; $h\nu\eu^{-h\nu/k_BT}$: congelado. (c) $u_\nu = (8\pi\nu^2/c^3)\langle E
\rangle$; Rayleigh--Jeans da $k_BT$ a cada modo. (d) Un modo no puede
excitarse con menos de un cuanto.
\end{solution}

\begin{solution}{exo:b2:boltzmann-factor:12}
(a) Solo $x = \mu B/k_BT$; $\ln 2$ y $0$. (b) $x = 0.67$: $S/Nk_B = 0.51$, luego
se retiran $0.18\,k_B$ por espín. (c) Con $B/T$ constante: \qty{10}{mK}. (d) El campo
propio de los espines (milésimas de tesla) sustituye a $B$ al final y limita el enfriamiento;
medir $M = n\mu\tanh(\mu B/k_BT)$ da $T$.
\end{solution}

\begin{solution}{pb:b2:boltzmann-factor:1}
\textbf{1.} $\dd P/\dd z = -nmg = -(mg/k_BT)P$; $H = \qty{8.4}{km}$.

\textbf{2.} $\eu^{-mgz/k_BT}$: la energía potencial de una molécula y la
temperatura del aire.

\textbf{3.} \qty{0.70}{}, \qty{0.52}{}, \qty{0.35}{bar}; $H\ln 2 = \qty{5.8}{km}$.

\textbf{4.} $P_0/g = \qty{1.03e4}{kg/m^2}$; $\times 4\pi R^2$: \qty{5.3e18}{kg}; $\rho_0H =
1.22 \times 8400 = \qty{1.03e4}{kg/m^2}$.

\textbf{5.} $5.3 \times 10^{18}/(29 \times 10^{-3}/N_A) = 1.1 \times 10^{44}$.

\textbf{6.} $H_{\text{O}_2} = 7.6$, $H_{\text{N}_2} = \qty{8.7}{km}$: la razón a \qty{50}{km}
sería $\eu^{-50/7.6 + 50/8.7} = 0.44$ de su valor en el suelo; la turbulencia mezcla
más deprisa de lo que separa la difusión.

\textbf{7.} Menor: el aire frío es más denso y la presión cae más deprisa.

\textbf{8.} \qty{2.5e25}{m^{-3}}; $\eu^{-11.9}$: \qty{1.7e20}{m^{-3}} --- no hay borde, solo un
convenio.

\textbf{9.} $11.2\sqrt{6370/6870} = \qty{10.8}{km/s}$.

\textbf{10.} Helio \qty{2040}{} y \qty{2300}{m/s}; nitrógeno \qty{770}{} y \qty{870}{m/s}.

\textbf{11.} $x = 5.3$: $4 \times 10^{-12}$; nitrógeno $\eu^{-196}$: ninguno.

\textbf{12.} $\tfrac14 \times 10^{12} \times 2300 \times 4 \times 10^{-12} \times \tfrac12 \approx
10^3$ átomos por metro cuadrado y segundo; para toda la Tierra y un año,
$2 \times 10^{25}$ átomos --- \qty{0.1}{kg}.

\textbf{13.} $5 \times 10^{-6} \times 1.1 \times 10^{44} = 5.5 \times 10^{38}$ átomos, \qty{3.7e12}{kg};
frente a \qty{0.1}{kg} al año, el tiempo de residencia sería absurdo.

\textbf{14.} El \omterm{def:b2:open-systems:cv}{régimen estacionario} exige una pérdida de \qty{3e6}{kg} al año: un tiempo
de residencia de cerca de un millón de años, que es el valor aceptado --- así que
la estimación térmica a \qty{1000}{K} se queda muy corta: la fuga la
dominan los episodios más calientes y los procesos no térmicos (iónicos).

\textbf{15.} $x = 3.75$: $\eu^{-14}$ en vez de $\eu^{-28}$ --- un factor $10^6$:
la pérdida es extraordinariamente sensible a la temperatura de la exosfera.

\textbf{16.} La Luna es caliente y ligera: todo escapa en tiempo
geológico; Titán es frío, con $x = 11$ para el nitrógeno.

\textbf{17.} $p_\pm = \eu^{\pm x}/2\cosh x$; $M = n\mu\tanh x$.

\textbf{18.} $x = \mu B/k_BT$: \qty{4.2}{A/m}, \qty{310}{A/m}, \qty{1.6e4}{A/m}
(saturación \qty{1.9e4}{A/m}); Curie para $x \ll 1$, es decir, $T \gg \qty{70}{mK}$.

\textbf{19.} $M$ depende solo de $T$ una vez conocido $B$; con la mayor pendiente cerca de
$x \approx 1$, $T \approx \mu B/k_B = \qty{70}{mK}$; se mide la
susceptibilidad con una bobina.

\textbf{20.} $\langle E\rangle = -N\mu B\tanh x$; $C = Nk_Bx^2/\cosh^2x$, con el máximo en
$x = 1.2$: \qty{56}{mK}.

\textbf{21.} $Nk_B\ln 2$ y $0$: todos los espines se alinean con el campo.

\textbf{22.} \qty{10}{mK}; el campo propio de los espines, de una milésima de tesla, sustituye a $B$
y fija el suelo.

\textbf{23.} $x = 1$ para $\mu_{\text{N}}B/k_B = \qty{0.37}{mK}$: Curie vale hasta un
milikelvin, y el termómetro es útil hasta los microkelvin.

\textbf{24.} $\Delta\nu/\nu = \sqrt{8k_BT\ln2/mc^2} \propto \sqrt T$: la anchura
de una raya mide la temperatura del gas.

\textbf{25.} Atmósfera, gas que escapa y espines: $\eu^{-E/k_BT}$ con $E = mgz$,
$\tfrac12mv^2$ y $\mp\mu B$.
\end{solution}
